\par\noindent\textbf{(1)} 利用BCH共轭展开或对 \(S^\dagger(r)aS(r)\) 关于 \(r\) 求微分，得到 \(S^\dagger aS=a\cosh r-a^\dagger\sinh r\)。同理 \(S^\dagger a^\dagger S=a^\dagger\cosh r-a\sinh r\)，两式保留正则对易关系。
\par\noindent\textbf{(2)} 定义 \(X=(a+a^\dagger)/\sqrt2\)、\(P=(a-a^\dagger)/(i\sqrt2)\)。Bogoliubov变换给 \(S^\dagger XS=e^{-r}X\)、\(S^\dagger PS=e^{r}P\)，所以压缩真空中 \((\Delta X)^2=e^{-2r}/2\)、\((\Delta P)^2=e^{2r}/2\)。
\par\noindent\textbf{(3)} 因为 \([X,P]=i\)，\(\Delta X\Delta P=1/2\) 恰饱和不确定度下界，但两个象限的涨落不等。平均粒子数 \(\langle a^\dagger a\rangle=\sinh^2r\)，压缩真空只含偶数光子数分量，这与相干态的Poisson光子统计不同。
\par\noindent\textbf{(4)} 若压缩参数是复数 \(\zeta=re^{i\phi}\)，最小噪声象限随 \(\phi/2\) 旋转。效率 \(\eta<1\) 的真空损耗通道把测量象限方差变成 \(V_{\rm out}=\eta e^{-2r}/2+(1-\eta)/2\)，因此实际压缩不能无限抵消损耗注入的真空噪声。
